% Generated deterministically from the frozen Markdown source.
% Responsible author: Carptopus.
\documentclass[10pt]{article}
\usepackage[a4paper,margin=24mm]{geometry}
\usepackage{microtype}
\usepackage{amsmath,amssymb,amsthm,mathtools}
\usepackage{needspace}
\usepackage{xcolor}
\usepackage{xurl}
\usepackage{fontspec}
\newfontfamily\cjkfont{Microsoft YaHei}
\usepackage[colorlinks=true,linkcolor=blue!55!black,citecolor=blue!55!black,urlcolor=blue!65!black]{hyperref}
\usepackage{fancyhdr}

\newenvironment{mdstatement}[1]
  {\par\Needspace{0.16\textheight}\noindent\textbf{#1.}\ \itshape}
  {\par\normalfont}

\setlength{\parindent}{0pt}
\setlength{\parskip}{0.42em}
\setlength{\emergencystretch}{6em}
\setlength{\headheight}{14pt}
\urlstyle{same}

\pagestyle{fancy}
\fancyhf{}
\fancyhead[L]{Carptopus}
\fancyhead[R]{Paving-matroid technical proof dossier}
\fancyfoot[C]{\thepage}

\title{Technical proof dossier for\\
\emph{Eventual real-rootedness in paving-matroid base polytopes}}
\author{Carptopus}
\date{7 October 2026}

\hypersetup{
  pdftitle={Technical proof dossier for Eventual real-rootedness in paving-matroid base polytopes},
  pdfauthor={Carptopus},
  pdfsubject={Frozen proof dependency dossier},
  pdfkeywords={paving matroid, Ehrhart h-star polynomial, proof dossier, SHA-256 closure}
}

\begin{document}
\maketitle
\begin{center}
\fcolorbox{black!35}{black!4}{\parbox{0.90\textwidth}{\centering\small
Integral technical supplement to the accompanying manuscript; not peer reviewed.}}
\end{center}


This dossier is part of the manuscript, not a historical bibliography.  The main
paper gives the definitions, exact theorem statements, and assembly arguments;
the long uniform contour estimates and endpoint kernel arguments are supplied by
the frozen proof records listed below.  Publication artifacts must contain this
file, every listed record, and the ordered 216-file SHA-256 manifest
\texttt{loops/\allowbreak{}PAVING-\allowbreak{}HSTAR-\allowbreak{}PUBLICATION-\allowbreak{}PREP-\allowbreak{}0001/\allowbreak{}A1-\allowbreak{}{\cjkfont 证\allowbreak{}明\allowbreak{}依\allowbreak{}赖\allowbreak{}闭\allowbreak{}包}.\allowbreak{}sha256}.
Changing any listed byte invalidates the manuscript audit.

The labels \texttt{PASS}, \texttt{candidate}, and \texttt{HOLD} appearing in the dated records are
historical workflow metadata.  They are not mathematical hypotheses and are not
used in the proofs.  Only the displayed identities, estimates, certificates,
and quantifiers are used.

\section{Exact Ehrhart correction formula}


\begin{itemize}
\item \texttt{loops/\allowbreak{}PAVING-\allowbreak{}HSTAR-\allowbreak{}HIGHER-\allowbreak{}RANK-\allowbreak{}CRITICAL-\allowbreak{}KERNEL-\allowbreak{}GATE-\allowbreak{}0001/\allowbreak{}S10-\allowbreak{}{\cjkfont 全\allowbreak{}秩\allowbreak{}内\allowbreak{}容\allowbreak{}矩\allowbreak{}阵\allowbreak{}的\allowbreak{}准\allowbreak{}确\allowbreak{}谱\allowbreak{}与\allowbreak{}复\allowbreak{}方\allowbreak{}差\allowbreak{}入\allowbreak{}口}.\allowbreak{}md}
\item its frozen revision \texttt{S10R1}
\item \texttt{S11-\allowbreak{}{\cjkfont 真\allowbreak{}实\allowbreak{}全\allowbreak{}容\allowbreak{}斥\allowbreak{}与\allowbreak{}低}Euler{\cjkfont 谱\allowbreak{}分\allowbreak{}离\allowbreak{}候\allowbreak{}选}.\allowbreak{}md}
\end{itemize}

These records prove the common-denominator formula (2.1)--(2.9), including the
full inclusion--exclusion over every effective hyperplane and every lower Euler
layer.  They are the algebraic input to all analytic estimates below.

\section{Fixed-rank central and bounded-excess estimates}


\begin{itemize}
\item \texttt{loops/\allowbreak{}PAVING-\allowbreak{}HSTAR-\allowbreak{}CROSS-\allowbreak{}RANK-\allowbreak{}FACT-\allowbreak{}GATE-\allowbreak{}0001/\allowbreak{}S18-\allowbreak{}{\cjkfont 中\allowbreak{}间\allowbreak{}小}t{\cjkfont 区\allowbreak{}的\allowbreak{}解\allowbreak{}析\allowbreak{}谱\allowbreak{}隙\allowbreak{}桥}.\allowbreak{}md}
\item \texttt{loops/\allowbreak{}PAVING-\allowbreak{}HSTAR-\allowbreak{}CROSS-\allowbreak{}RANK-\allowbreak{}FACT-\allowbreak{}GATE-\allowbreak{}0001/\allowbreak{}S14-\allowbreak{}{\cjkfont 增\allowbreak{}长\allowbreak{}秩\allowbreak{}紧\allowbreak{}中\allowbreak{}带\allowbreak{}的\allowbreak{}定\allowbreak{}量\allowbreak{}谱\allowbreak{}隙\allowbreak{}与\allowbreak{}真\allowbreak{}实\allowbreak{}保\allowbreak{}号}.\allowbreak{}md}
\item \texttt{loops/\allowbreak{}PAVING-\allowbreak{}HSTAR-\allowbreak{}CROSS-\allowbreak{}RANK-\allowbreak{}FACT-\allowbreak{}GATE-\allowbreak{}0001/\allowbreak{}S15-\allowbreak{}{\cjkfont 增\allowbreak{}长\allowbreak{}秩\allowbreak{}倒\allowbreak{}端\allowbreak{}移\allowbreak{}动\allowbreak{}保\allowbreak{}号\allowbreak{}的\allowbreak{}双\allowbreak{}区\allowbreak{}定\allowbreak{}量\allowbreak{}桥}.\allowbreak{}md}
\item \texttt{loops/\allowbreak{}PAVING-\allowbreak{}HSTAR-\allowbreak{}HIGHER-\allowbreak{}RANK-\allowbreak{}CRITICAL-\allowbreak{}KERNEL-\allowbreak{}GATE-\allowbreak{}0001/\allowbreak{}S90-\allowbreak{}{\cjkfont 全\allowbreak{}部\allowbreak{}固\allowbreak{}定\allowbreak{}秩\allowbreak{}全\allowbreak{}缺\allowbreak{}口\allowbreak{}临\allowbreak{}界\allowbreak{}核\allowbreak{}联\allowbreak{}合\allowbreak{}候\allowbreak{}选}.\allowbreak{}md}
\item \texttt{loops/\allowbreak{}PAVING-\allowbreak{}HSTAR-\allowbreak{}HIGHER-\allowbreak{}RANK-\allowbreak{}CRITICAL-\allowbreak{}KERNEL-\allowbreak{}GATE-\allowbreak{}0001/\allowbreak{}S59-\allowbreak{}{\cjkfont 全\allowbreak{}固\allowbreak{}定\allowbreak{}秩\allowbreak{}非\allowbreak{}负\allowbreak{}临\allowbreak{}界\allowbreak{}超\allowbreak{}额\allowbreak{}带\allowbreak{}实\allowbreak{}际\allowbreak{}全\allowbreak{}轴\allowbreak{}候\allowbreak{}选}.\allowbreak{}md}
\item \texttt{loops/\allowbreak{}PAVING-\allowbreak{}HSTAR-\allowbreak{}HIGHER-\allowbreak{}RANK-\allowbreak{}CRITICAL-\allowbreak{}KERNEL-\allowbreak{}GATE-\allowbreak{}0001/\allowbreak{}S91-\allowbreak{}{\cjkfont 全\allowbreak{}固\allowbreak{}定\allowbreak{}秩\allowbreak{}有\allowbreak{}界\allowbreak{}临\allowbreak{}界\allowbreak{}超\allowbreak{}额\allowbreak{}的\allowbreak{}真\allowbreak{}实\allowbreak{}全\allowbreak{}轴\allowbreak{}联\allowbreak{}合\allowbreak{}候\allowbreak{}选}.\allowbreak{}md}
\end{itemize}

The S91 proof is independent of Proposition 3.1 outside its stated density
range.  In the band $D_u\le s\le D_u+C$, it uses the exact fixed increment
parameters

\[
 b_u=r\lceil n/r\rceil-n,
 \quad \ell=\left\lceil\frac{\delta+1-b_u}{r-1}\right\rceil,
 \quad c=\lceil n/r\rceil+\ell,
 \quad b=b_u+r\ell,
 \quad a=(r-1)\ell+b_u-\delta-1,
\]

and proves directly that the bounded shift of the highest cutoff changes the
two-moving local expansion by $\exp(O(\sigma))-1=o(1)$, while the lower-Euler
and packing errors remain (o(1/r)).  This is the input used in Lemma 5.1.

\section{Fixed complement and rank-three interfaces}


\begin{itemize}
\item \texttt{loops/\allowbreak{}PAVING-\allowbreak{}HSTAR-\allowbreak{}FIXED-\allowbreak{}OUTSIDE-\allowbreak{}GATE-\allowbreak{}0001/\allowbreak{}S1-\allowbreak{}{\cjkfont 固\allowbreak{}定\allowbreak{}外\allowbreak{}部\allowbreak{}数\allowbreak{}二\allowbreak{}项\allowbreak{}式\allowbreak{}矩\allowbreak{}与\allowbreak{}完\allowbreak{}整\allowbreak{}候\allowbreak{}选}.\allowbreak{}md}
\item \texttt{loops/\allowbreak{}PAVING-\allowbreak{}HSTAR-\allowbreak{}FIXED-\allowbreak{}OUTSIDE-\allowbreak{}GATE-\allowbreak{}0001/\allowbreak{}S2-\allowbreak{}{\cjkfont 秩\allowbreak{}三\allowbreak{}真\allowbreak{}实\allowbreak{}校\allowbreak{}正\allowbreak{}与\allowbreak{}两\allowbreak{}端\allowbreak{}极\allowbreak{}限\allowbreak{}候\allowbreak{}选}.\allowbreak{}md}
\item \texttt{loops/\allowbreak{}PAVING-\allowbreak{}HSTAR-\allowbreak{}MOVING-\allowbreak{}CRITICAL-\allowbreak{}GATE-\allowbreak{}0001/\allowbreak{}S22-\allowbreak{}{\cjkfont 秩\allowbreak{}三\allowbreak{}最\allowbreak{}终\allowbreak{}全\allowbreak{}类\allowbreak{}的\allowbreak{}联\allowbreak{}合\allowbreak{}候\allowbreak{}选\allowbreak{}与\allowbreak{}证\allowbreak{}据\allowbreak{}门}.\allowbreak{}md}
\item its mathematical revision \texttt{S22R1-\allowbreak{}{\cjkfont 巨}H{\cjkfont 正\allowbreak{}因\allowbreak{}子\allowbreak{}依\allowbreak{}赖\allowbreak{}绑\allowbreak{}定}.\allowbreak{}md}
\item \texttt{S24-\allowbreak{}{\cjkfont 最\allowbreak{}终\allowbreak{}秩\allowbreak{}三\allowbreak{}自\allowbreak{}然\allowbreak{}整\allowbreak{}类\allowbreak{}结\allowbreak{}算\allowbreak{}与\allowbreak{}跨\allowbreak{}秩\allowbreak{}接\allowbreak{}口}.\allowbreak{}md}
\item \texttt{loops/\allowbreak{}PAVING-\allowbreak{}HSTAR-\allowbreak{}GIANT-\allowbreak{}H-\allowbreak{}DEPENDENCIES-\allowbreak{}GATE-\allowbreak{}0001/\allowbreak{}S9-\allowbreak{}{\cjkfont 全\allowbreak{}部\allowbreak{}亚\allowbreak{}线\allowbreak{}性\allowbreak{}整\allowbreak{}类\allowbreak{}与\allowbreak{}近\allowbreak{}一\allowbreak{}宏\allowbreak{}观\allowbreak{}范\allowbreak{}围\allowbreak{}全\allowbreak{}轴\allowbreak{}候\allowbreak{}选}.\allowbreak{}md}
\item its dependency revision \texttt{S9R1-\allowbreak{}{\cjkfont 全\allowbreak{}部\allowbreak{}亚\allowbreak{}线\allowbreak{}性\allowbreak{}完\allowbreak{}整\allowbreak{}包\allowbreak{}依\allowbreak{}赖\allowbreak{}修\allowbreak{}订}.\allowbreak{}md}
\item \texttt{loops/\allowbreak{}PAVING-\allowbreak{}HSTAR-\allowbreak{}GIANT-\allowbreak{}H-\allowbreak{}DEPENDENCIES-\allowbreak{}GATE-\allowbreak{}0001/\allowbreak{}S16-\allowbreak{}{\cjkfont 临\allowbreak{}界\allowbreak{}邻\allowbreak{}域\allowbreak{}整\allowbreak{}类\allowbreak{}全\allowbreak{}轴\allowbreak{}同\allowbreak{}伦\allowbreak{}候\allowbreak{}选}.\allowbreak{}md}
\item \texttt{loops/\allowbreak{}PAVING-\allowbreak{}HSTAR-\allowbreak{}GIANT-\allowbreak{}H-\allowbreak{}DEPENDENCIES-\allowbreak{}GATE-\allowbreak{}0001/\allowbreak{}S19-\allowbreak{}{\cjkfont 全\allowbreak{}正\allowbreak{}比\allowbreak{}例\allowbreak{}双\allowbreak{}端\allowbreak{}点\allowbreak{}间\allowbreak{}隙\allowbreak{}完\allowbreak{}整\allowbreak{}候\allowbreak{}选}.\allowbreak{}md}
\item its mathematical revision \texttt{S19R1-\allowbreak{}{\cjkfont 大\allowbreak{}域\allowbreak{}正\allowbreak{}因\allowbreak{}子\allowbreak{}与\allowbreak{}双\allowbreak{}端\allowbreak{}点\allowbreak{}证\allowbreak{}书\allowbreak{}修\allowbreak{}订}.\allowbreak{}md}
\item \texttt{loops/\allowbreak{}PAVING-\allowbreak{}HSTAR-\allowbreak{}GIANT-\allowbreak{}H-\allowbreak{}DEPENDENCIES-\allowbreak{}GATE-\allowbreak{}0001/\allowbreak{}S20-\allowbreak{}{\cjkfont 全\allowbreak{}部\allowbreak{}巨}H{\cjkfont 依\allowbreak{}赖\allowbreak{}整\allowbreak{}类\allowbreak{}统\allowbreak{}一\allowbreak{}实\allowbreak{}际\allowbreak{}全\allowbreak{}轴\allowbreak{}候\allowbreak{}选}.\allowbreak{}md}
\end{itemize}

For each fixed outside size $q\ge2$, S1 proves the rank-$r\ge4$ theorem by
an exact finite binomial-moment expansion and a three-region homotopy.  S2 gives
the separate rank-three proof; it does not invoke the false (k=2) general
perturbation box.  S20 supplies the all-attached-hyperplanes giant-region
theorem.  S22 supplies the remaining rank-three failure-sequence exhaustion,
including the continuous negative-axis certificate and the exact complex-
circle counts.

\section{Giant-hyperplane middle and endpoint estimates}


\begin{itemize}
\item \texttt{loops/\allowbreak{}PAVING-\allowbreak{}HSTAR-\allowbreak{}TWO-\allowbreak{}MACRO-\allowbreak{}GATE-\allowbreak{}0001/\allowbreak{}S3-\allowbreak{}{\cjkfont 秩\allowbreak{}至\allowbreak{}少\allowbreak{}四\allowbreak{}全\allowbreak{}部\allowbreak{}宏\allowbreak{}观\allowbreak{}谱\allowbreak{}统\allowbreak{}一\allowbreak{}候\allowbreak{}选}.\allowbreak{}md}
\item \texttt{loops/\allowbreak{}PAVING-\allowbreak{}HSTAR-\allowbreak{}GENERAL-\allowbreak{}RANK-\allowbreak{}GIANT-\allowbreak{}GEOMETRY-\allowbreak{}GATE-\allowbreak{}0001/\allowbreak{}S1-\allowbreak{}{\cjkfont 一\allowbreak{}般\allowbreak{}秩\allowbreak{}全\allowbreak{}圆\allowbreak{}最\allowbreak{}小\allowbreak{}模\allowbreak{}的\allowbreak{}角\allowbreak{}区\allowbreak{}排\allowbreak{}除\allowbreak{}候\allowbreak{}选}.\allowbreak{}md}
\item \texttt{loops/\allowbreak{}PAVING-\allowbreak{}HSTAR-\allowbreak{}GENERAL-\allowbreak{}RANK-\allowbreak{}GIANT-\allowbreak{}GEOMETRY-\allowbreak{}GATE-\allowbreak{}0001/\allowbreak{}S2-\allowbreak{}{\cjkfont 一\allowbreak{}般\allowbreak{}秩\allowbreak{}准\allowbreak{}确\allowbreak{}累\allowbreak{}计\allowbreak{}积\allowbreak{}分\allowbreak{}与\allowbreak{}移\allowbreak{}动\allowbreak{}中\allowbreak{}央\allowbreak{}预\allowbreak{}算\allowbreak{}候\allowbreak{}选}.\allowbreak{}md}
\item \texttt{loops/\allowbreak{}PAVING-\allowbreak{}HSTAR-\allowbreak{}GENERAL-\allowbreak{}RANK-\allowbreak{}GIANT-\allowbreak{}MOVING-\allowbreak{}GATE-\allowbreak{}0001/\allowbreak{}S3-\allowbreak{}{\cjkfont 一\allowbreak{}般\allowbreak{}秩\allowbreak{}小\allowbreak{}端\allowbreak{}分\allowbreak{}离\allowbreak{}累\allowbreak{}计\allowbreak{}完\allowbreak{}整\allowbreak{}积\allowbreak{}分\allowbreak{}候\allowbreak{}选}.\allowbreak{}md}
\item \texttt{loops/\allowbreak{}PAVING-\allowbreak{}HSTAR-\allowbreak{}GENERAL-\allowbreak{}RANK-\allowbreak{}GIANT-\allowbreak{}MOVING-\allowbreak{}GATE-\allowbreak{}0001/\allowbreak{}S4-\allowbreak{}{\cjkfont 一\allowbreak{}般\allowbreak{}秩\allowbreak{}正\allowbreak{}紧\allowbreak{}驻\allowbreak{}点\allowbreak{}模\allowbreak{}完\allowbreak{}整\allowbreak{}积\allowbreak{}分\allowbreak{}候\allowbreak{}选}.\allowbreak{}md}
\item \texttt{loops/\allowbreak{}PAVING-\allowbreak{}HSTAR-\allowbreak{}GENERAL-\allowbreak{}RANK-\allowbreak{}GIANT-\allowbreak{}MOVING-\allowbreak{}GATE-\allowbreak{}0001/\allowbreak{}S5-\allowbreak{}{\cjkfont 一\allowbreak{}般\allowbreak{}秩\allowbreak{}大\allowbreak{}端\allowbreak{}有\allowbreak{}界}K{\cjkfont 的\allowbreak{}完\allowbreak{}整\allowbreak{}积\allowbreak{}分\allowbreak{}候\allowbreak{}选}.\allowbreak{}md}
\item \texttt{loops/\allowbreak{}PAVING-\allowbreak{}HSTAR-\allowbreak{}GENERAL-\allowbreak{}RANK-\allowbreak{}GIANT-\allowbreak{}MOVING-\allowbreak{}GATE-\allowbreak{}0001/\allowbreak{}S6-\allowbreak{}{\cjkfont 大\allowbreak{}端\allowbreak{}共\allowbreak{}同\allowbreak{}无\allowbreak{}界\allowbreak{}多\allowbreak{}峰\allowbreak{}相\allowbreak{}对\allowbreak{}尾\allowbreak{}候\allowbreak{}选}.\allowbreak{}md}
\item \texttt{loops/\allowbreak{}PAVING-\allowbreak{}HSTAR-\allowbreak{}GENERAL-\allowbreak{}RANK-\allowbreak{}GIANT-\allowbreak{}MOVING-\allowbreak{}GATE-\allowbreak{}0001/\allowbreak{}S7-\allowbreak{}{\cjkfont 一\allowbreak{}般\allowbreak{}秩\allowbreak{}有\allowbreak{}界\allowbreak{}累\allowbreak{}计\allowbreak{}窗\allowbreak{}与\allowbreak{}分\allowbreak{}离\allowbreak{}域\allowbreak{}接\allowbreak{}合\allowbreak{}候\allowbreak{}选}.\allowbreak{}md}
\item \texttt{loops/\allowbreak{}PAVING-\allowbreak{}HSTAR-\allowbreak{}GENERAL-\allowbreak{}RANK-\allowbreak{}GIANT-\allowbreak{}MOVING-\allowbreak{}GATE-\allowbreak{}0001/\allowbreak{}S8-\allowbreak{}{\cjkfont 全\allowbreak{}一\allowbreak{}般\allowbreak{}秩\allowbreak{}分\allowbreak{}离\allowbreak{}累\allowbreak{}计\allowbreak{}的\allowbreak{}联\allowbreak{}合\allowbreak{}量\allowbreak{}词\allowbreak{}候\allowbreak{}选}.\allowbreak{}md}
\item \texttt{loops/\allowbreak{}PAVING-\allowbreak{}HSTAR-\allowbreak{}GENERAL-\allowbreak{}RANK-\allowbreak{}GIANT-\allowbreak{}MOVING-\allowbreak{}GATE-\allowbreak{}0001/\allowbreak{}S9-\allowbreak{}{\cjkfont 最\allowbreak{}高\allowbreak{}累\allowbreak{}计\allowbreak{}全\allowbreak{}移\allowbreak{}动\allowbreak{}两\allowbreak{}域\allowbreak{}接\allowbreak{}合\allowbreak{}候\allowbreak{}选}.\allowbreak{}md}
\item \texttt{loops/\allowbreak{}PAVING-\allowbreak{}HSTAR-\allowbreak{}GENERAL-\allowbreak{}RANK-\allowbreak{}GIANT-\allowbreak{}ACTUAL-\allowbreak{}GATE-\allowbreak{}0001/\allowbreak{}S6-\allowbreak{}{\cjkfont 圆\allowbreak{}盘\allowbreak{}极\allowbreak{}化\allowbreak{}直\allowbreak{}接\allowbreak{}排\allowbreak{}除\allowbreak{}多\allowbreak{}标\allowbreak{}记\allowbreak{}谱\allowbreak{}候\allowbreak{}选}.\allowbreak{}md}
\item \texttt{loops/\allowbreak{}PAVING-\allowbreak{}HSTAR-\allowbreak{}GENERAL-\allowbreak{}RANK-\allowbreak{}GIANT-\allowbreak{}ACTUAL-\allowbreak{}GATE-\allowbreak{}0001/\allowbreak{}S7-\allowbreak{}{\cjkfont 全\allowbreak{}部\allowbreak{}实\allowbreak{}际\allowbreak{}移\allowbreak{}动\allowbreak{}两\allowbreak{}域\allowbreak{}联\allowbreak{}合\allowbreak{}候\allowbreak{}选}.\allowbreak{}md}
\item \texttt{loops/\allowbreak{}PAVING-\allowbreak{}HSTAR-\allowbreak{}GENERAL-\allowbreak{}RANK-\allowbreak{}GIANT-\allowbreak{}ACTUAL-\allowbreak{}GATE-\allowbreak{}0001/\allowbreak{}S8-\allowbreak{}{\cjkfont 实\allowbreak{}际\allowbreak{}移\allowbreak{}动\allowbreak{}中\allowbreak{}带\allowbreak{}结\allowbreak{}算\allowbreak{}与\allowbreak{}端\allowbreak{}层\allowbreak{}入\allowbreak{}口}.\allowbreak{}md}
\item \texttt{loops/\allowbreak{}PAVING-\allowbreak{}HSTAR-\allowbreak{}GENERAL-\allowbreak{}RANK-\allowbreak{}GIANT-\allowbreak{}ENDPOINT-\allowbreak{}GATE-\allowbreak{}0001/\allowbreak{}S2-\allowbreak{}{\cjkfont 一\allowbreak{}般}Wright{\cjkfont 相\allowbreak{}位\allowbreak{}圆\allowbreak{}准\allowbreak{}确\allowbreak{}索\allowbreak{}引\allowbreak{}候\allowbreak{}选}.\allowbreak{}md}
\item \texttt{loops/\allowbreak{}PAVING-\allowbreak{}HSTAR-\allowbreak{}GENERAL-\allowbreak{}RANK-\allowbreak{}GIANT-\allowbreak{}ENDPOINT-\allowbreak{}GATE-\allowbreak{}0001/\allowbreak{}S4-\allowbreak{}{\cjkfont 一\allowbreak{}般\allowbreak{}秩\allowbreak{}发\allowbreak{}散\allowbreak{}巨}H{\cjkfont 倒\allowbreak{}端\allowbreak{}与\allowbreak{}准\allowbreak{}确\allowbreak{}圆\allowbreak{}候\allowbreak{}选}.\allowbreak{}md}
\item \texttt{loops/\allowbreak{}PAVING-\allowbreak{}HSTAR-\allowbreak{}GENERAL-\allowbreak{}RANK-\allowbreak{}GIANT-\allowbreak{}ENDPOINT-\allowbreak{}GATE-\allowbreak{}0001/\allowbreak{}S5-\allowbreak{}{\cjkfont 一\allowbreak{}般\allowbreak{}秩\allowbreak{}巨}H{\cjkfont 双\allowbreak{}端\allowbreak{}连\allowbreak{}续\allowbreak{}相\allowbreak{}位\allowbreak{}接\allowbreak{}合\allowbreak{}候\allowbreak{}选}.\allowbreak{}md}
\item \texttt{loops/\allowbreak{}PAVING-\allowbreak{}HSTAR-\allowbreak{}GENERAL-\allowbreak{}RANK-\allowbreak{}GIANT-\allowbreak{}ENDPOINT-\allowbreak{}GATE-\allowbreak{}0001/\allowbreak{}S6-\allowbreak{}{\cjkfont 一\allowbreak{}般\allowbreak{}秩\allowbreak{}亚\allowbreak{}临\allowbreak{}界\allowbreak{}巨}H{\cjkfont 实\allowbreak{}际\allowbreak{}全\allowbreak{}轴\allowbreak{}联\allowbreak{}合\allowbreak{}候\allowbreak{}选}.\allowbreak{}md}
\item \texttt{loops/\allowbreak{}PAVING-\allowbreak{}HSTAR-\allowbreak{}GENERAL-\allowbreak{}RANK-\allowbreak{}GIANT-\allowbreak{}CRITICAL-\allowbreak{}MOVING-\allowbreak{}GATE-\allowbreak{}0001/\allowbreak{}S1-\allowbreak{}{\cjkfont 全\allowbreak{}秩\allowbreak{}临\allowbreak{}界\allowbreak{}与\allowbreak{}分\allowbreak{}离\allowbreak{}基\allowbreak{}础\allowbreak{}相\allowbreak{}位\allowbreak{}符\allowbreak{}号\allowbreak{}候\allowbreak{}选}.\allowbreak{}md}
\item \texttt{loops/\allowbreak{}PAVING-\allowbreak{}HSTAR-\allowbreak{}GENERAL-\allowbreak{}RANK-\allowbreak{}GIANT-\allowbreak{}CRITICAL-\allowbreak{}MOVING-\allowbreak{}GATE-\allowbreak{}0001/\allowbreak{}S2-\allowbreak{}{\cjkfont 全\allowbreak{}秩\allowbreak{}巨}H{\cjkfont 临\allowbreak{}界\allowbreak{}宏\allowbreak{}观\allowbreak{}基\allowbreak{}础\allowbreak{}圈\allowbreak{}准\allowbreak{}确\allowbreak{}索\allowbreak{}引\allowbreak{}候\allowbreak{}选}.\allowbreak{}md}
\item \texttt{loops/\allowbreak{}PAVING-\allowbreak{}HSTAR-\allowbreak{}GENERAL-\allowbreak{}RANK-\allowbreak{}GIANT-\allowbreak{}CRITICAL-\allowbreak{}MOVING-\allowbreak{}GATE-\allowbreak{}0001/\allowbreak{}S3-\allowbreak{}{\cjkfont 全\allowbreak{}秩\allowbreak{}慢\allowbreak{}发\allowbreak{}散\allowbreak{}临\allowbreak{}界\allowbreak{}倒\allowbreak{}端\allowbreak{}准\allowbreak{}确\allowbreak{}相\allowbreak{}位\allowbreak{}候\allowbreak{}选}.\allowbreak{}md}
\item \texttt{loops/\allowbreak{}PAVING-\allowbreak{}HSTAR-\allowbreak{}GENERAL-\allowbreak{}RANK-\allowbreak{}GIANT-\allowbreak{}CRITICAL-\allowbreak{}MOVING-\allowbreak{}GATE-\allowbreak{}0001/\allowbreak{}S4-\allowbreak{}{\cjkfont 全\allowbreak{}秩\allowbreak{}临\allowbreak{}界\allowbreak{}移\allowbreak{}动\allowbreak{}实\allowbreak{}际\allowbreak{}全\allowbreak{}轴\allowbreak{}联\allowbreak{}合\allowbreak{}候\allowbreak{}选}.\allowbreak{}md}
\end{itemize}

Together these records prove the four-parameter finite gate in Proposition 3.4,
the exact Wright/Beta direct circles, the reciprocal factorial circles, and the
continuity/first-hit assembly.  In particular, the endpoint counts are full
complex-circle argument-principle counts, not deductions from a single
negative-axis sign.

\section{Rank-four exhaustion}


\begin{itemize}
\item \texttt{loops/\allowbreak{}PAVING-\allowbreak{}HSTAR-\allowbreak{}HIGHER-\allowbreak{}RANK-\allowbreak{}CRITICAL-\allowbreak{}KERNEL-\allowbreak{}GATE-\allowbreak{}0001/\allowbreak{}S50-\allowbreak{}{\cjkfont 稀\allowbreak{}外\allowbreak{}部\allowbreak{}原\allowbreak{}端}B{\cjkfont 核\allowbreak{}的\allowbreak{}准\allowbreak{}确\allowbreak{}三\allowbreak{}支\allowbreak{}相\allowbreak{}位\allowbreak{}索\allowbreak{}引\allowbreak{}候\allowbreak{}选}.\allowbreak{}md}
\item \texttt{loops/\allowbreak{}PAVING-\allowbreak{}HSTAR-\allowbreak{}HIGHER-\allowbreak{}RANK-\allowbreak{}CRITICAL-\allowbreak{}KERNEL-\allowbreak{}GATE-\allowbreak{}0001/\allowbreak{}S51-\allowbreak{}Wright{\cjkfont 原\allowbreak{}端\allowbreak{}的\allowbreak{}全\allowbreak{}扇\allowbreak{}区\allowbreak{}无\allowbreak{}穷\allowbreak{}相\allowbreak{}位\allowbreak{}候\allowbreak{}选}.\allowbreak{}md}
\item \texttt{loops/\allowbreak{}PAVING-\allowbreak{}HSTAR-\allowbreak{}HIGHER-\allowbreak{}RANK-\allowbreak{}CRITICAL-\allowbreak{}KERNEL-\allowbreak{}GATE-\allowbreak{}0001/\allowbreak{}S52-\allowbreak{}{\cjkfont 稀\allowbreak{}外\allowbreak{}部\allowbreak{}实\allowbreak{}际\allowbreak{}全\allowbreak{}轴\allowbreak{}的\allowbreak{}准\allowbreak{}确\allowbreak{}双\allowbreak{}端\allowbreak{}计\allowbreak{}数\allowbreak{}候\allowbreak{}选}.\allowbreak{}md}
\item \texttt{loops/\allowbreak{}PAVING-\allowbreak{}HSTAR-\allowbreak{}HIGHER-\allowbreak{}RANK-\allowbreak{}CRITICAL-\allowbreak{}KERNEL-\allowbreak{}GATE-\allowbreak{}0001/\allowbreak{}S53-\allowbreak{}{\cjkfont 巨}H{\cjkfont 固\allowbreak{}定\allowbreak{}宏\allowbreak{}观\allowbreak{}比\allowbreak{}例\allowbreak{}的\allowbreak{}原\allowbreak{}端\allowbreak{}索\allowbreak{}引\allowbreak{}与\allowbreak{}实\allowbreak{}际\allowbreak{}全\allowbreak{}轴\allowbreak{}候\allowbreak{}选}.\allowbreak{}md}
\item \texttt{loops/\allowbreak{}PAVING-\allowbreak{}HSTAR-\allowbreak{}HIGHER-\allowbreak{}RANK-\allowbreak{}CRITICAL-\allowbreak{}KERNEL-\allowbreak{}GATE-\allowbreak{}0001/\allowbreak{}S54-\allowbreak{}{\cjkfont 临\allowbreak{}界\allowbreak{}与\allowbreak{}分\allowbreak{}离\allowbreak{}共\allowbreak{}同\allowbreak{}基\allowbreak{}础\allowbreak{}相\allowbreak{}位\allowbreak{}符\allowbreak{}号\allowbreak{}门\allowbreak{}候\allowbreak{}选}.\allowbreak{}md}
\item \texttt{loops/\allowbreak{}PAVING-\allowbreak{}HSTAR-\allowbreak{}HIGHER-\allowbreak{}RANK-\allowbreak{}CRITICAL-\allowbreak{}KERNEL-\allowbreak{}GATE-\allowbreak{}0001/\allowbreak{}S55-\allowbreak{}{\cjkfont 全\allowbreak{}部\allowbreak{}巨}H{\cjkfont 宏\allowbreak{}观\allowbreak{}比\allowbreak{}例\allowbreak{}的\allowbreak{}基\allowbreak{}础\allowbreak{}圆\allowbreak{}准\allowbreak{}确\allowbreak{}索\allowbreak{}引\allowbreak{}候\allowbreak{}选}.\allowbreak{}md}
\item \texttt{loops/\allowbreak{}PAVING-\allowbreak{}HSTAR-\allowbreak{}HIGHER-\allowbreak{}RANK-\allowbreak{}CRITICAL-\allowbreak{}KERNEL-\allowbreak{}GATE-\allowbreak{}0001/\allowbreak{}S56-\allowbreak{}{\cjkfont 移\allowbreak{}动\allowbreak{}临\allowbreak{}界\allowbreak{}巨}H{\cjkfont 的\allowbreak{}准\allowbreak{}确\allowbreak{}全\allowbreak{}轴\allowbreak{}计\allowbreak{}数\allowbreak{}候\allowbreak{}选}.\allowbreak{}md}
\item \texttt{loops/\allowbreak{}PAVING-\allowbreak{}HSTAR-\allowbreak{}HIGHER-\allowbreak{}RANK-\allowbreak{}CRITICAL-\allowbreak{}KERNEL-\allowbreak{}GATE-\allowbreak{}0001/\allowbreak{}S57-\allowbreak{}{\cjkfont 完\allowbreak{}整\allowbreak{}秩\allowbreak{}四\allowbreak{}自\allowbreak{}然\allowbreak{}整\allowbreak{}类\allowbreak{}的\allowbreak{}联\allowbreak{}合\allowbreak{}穷\allowbreak{}尽\allowbreak{}候\allowbreak{}选}.\allowbreak{}md}
\end{itemize}

The S50--S56 records prove, respectively, the sparse-outside direct endpoint,
Wright-sector phase, sparse-outside full axis, fixed macroscopic profile,
critical/separated sign gate, macroscopic endpoint index, and moving-critical
full-axis theorem.  S57 gives the exhaustive subsequence argument and the
disconnected classification used in Theorem 1.1.

\section{Growing-rank window}


\begin{itemize}
\item \texttt{loops/\allowbreak{}PAVING-\allowbreak{}HSTAR-\allowbreak{}CROSS-\allowbreak{}RANK-\allowbreak{}FACT-\allowbreak{}GATE-\allowbreak{}0001/\allowbreak{}S4-\allowbreak{}{\cjkfont 增\allowbreak{}长\allowbreak{}秩\allowbreak{}真\allowbreak{}实\allowbreak{}倒\allowbreak{}端\allowbreak{}定\allowbreak{}量\allowbreak{}桥\allowbreak{}候\allowbreak{}选}.\allowbreak{}md}
\item \texttt{loops/\allowbreak{}PAVING-\allowbreak{}HSTAR-\allowbreak{}CROSS-\allowbreak{}RANK-\allowbreak{}FACT-\allowbreak{}GATE-\allowbreak{}0001/\allowbreak{}S15-\allowbreak{}{\cjkfont 增\allowbreak{}长\allowbreak{}秩\allowbreak{}倒\allowbreak{}端\allowbreak{}移\allowbreak{}动\allowbreak{}保\allowbreak{}号\allowbreak{}的\allowbreak{}双\allowbreak{}区\allowbreak{}定\allowbreak{}量\allowbreak{}桥}.\allowbreak{}md}
\item \texttt{loops/\allowbreak{}PAVING-\allowbreak{}HSTAR-\allowbreak{}CROSS-\allowbreak{}RANK-\allowbreak{}FACT-\allowbreak{}GATE-\allowbreak{}0001/\allowbreak{}S19-\allowbreak{}{\cjkfont 显\allowbreak{}式\allowbreak{}增\allowbreak{}长\allowbreak{}秩\allowbreak{}下\allowbreak{}的\allowbreak{}渐\allowbreak{}近\allowbreak{}全\allowbreak{}实\allowbreak{}根\allowbreak{}候\allowbreak{}选}.\allowbreak{}md}
\item \texttt{loops/\allowbreak{}PAVING-\allowbreak{}HSTAR-\allowbreak{}CROSS-\allowbreak{}RANK-\allowbreak{}FACT-\allowbreak{}GATE-\allowbreak{}0001/\allowbreak{}S20-\allowbreak{}S17{\cjkfont 门\allowbreak{}下\allowbreak{}的\allowbreak{}完\allowbreak{}整\allowbreak{}原\allowbreak{}端\allowbreak{}保\allowbreak{}号}.\allowbreak{}md}
\end{itemize}

S19 gives the complete arithmetic assembly: the direct grid contributes
(D-K-J) roots, the reciprocal disk contributes (K-1), and the last sign
gap contributes one.  The reciprocal estimate is a uniform Rouché inequality
on the whole circle, as recorded in S4; no finite negative-axis test substitutes
for that estimate.

\section{Finite exact guards}


The manifest contains six frozen certificate programs and their exact outputs.  They
certify only the rational inequalities and continuous box covers stated in
their headers.  Each includes its prescribed positive, negative, or destructive
control.  The analytic deductions from those inequalities remain in D1--D6.

\end{document}
